    \section{Собственные векторы и характеристическое уравнение}

    \begin{definition}[Собственное значение и собственный вектор]
Пусть $A \in \mathbb{R}^{n \times n}$ — квадратная матрица.
Ненулевой вектор $v \in \mathbb{C}^n$ ($v \ne 0$) называется \textbf{собственным вектором} матрицы $A$, а число $\lambda
\in \mathbb{C}$ — соответствующим ему \textbf{собственным значением}, если выполнено соотношение:
\[ A v = \lambda v. \]
    \end{definition}

    \begin{theorem}[Вывод характеристического уравнения]
Равенство $Av = \lambda v$ эквивалентно однородной системе:
\[ (A - \lambda I_n) v = 0. \]
По определению собственного вектора требуется найти \textbf{ненулевое} решение $v \ne 0$.
Ненулевое решение однородной системы существует тогда и только тогда, когда ядро матрицы $(A - \lambda I_n)$
нетривиально:
\[ \ker(A - \lambda I_n) \ne \{0\} \iff \rank(A - \lambda I_n) < n. \]
В силу критерия вырожденности квадратной матрицы это эквивалентно равенству нулю ее определителя:
\[ p_A(\lambda) \coloneqq \det(A - \lambda I_n) = 0. \]
Уравнение $\det(A - \lambda I_n) = 0$ называется \textbf{характеристическим уравнением} матрицы $A$, а $p_A(\lambda)$ —
\textbf{характеристическим многочленом} степени $n$.
    \end{theorem}

    \section{Алгебраическая и геометрическая кратности}

    \begin{definition}[Кратности собственных значений]
Пусть $\lambda_0$ — собственное значение матрицы $A$.
\begin{enumerate}
    \item \textbf{Алгебраическая кратность} $\mu_A(\lambda_0)$ (или $\operatorname{alg\_mult}(\lambda_0)$) — это
    кратность корня $\lambda_0$ в характеристическом многочлене $p_A(\lambda)$:
    \[ p_A(\lambda) = (\lambda - \lambda_0)^{\mu_A(\lambda_0)} q(\lambda), \quad q(\lambda_0) \ne 0. \]
    \item \textbf{Геометрическая кратность} $\gamma_A(\lambda_0)$ (или $\operatorname{geom\_mult}(\lambda_0)$) — это
    размерность соответствующего собственного подпространства $E_{\lambda_0} = \ker(A - \lambda_0 I_n)$:
    \[ \gamma_A(\lambda_0) \coloneqq \dim \ker(A - \lambda_0 I_n) = n - \rank(A - \lambda_0 I_n). \]
\end{enumerate}
    \end{definition}

    \begin{theorem}[Неравенство кратностей]
Для любого собственного значения $\lambda_0$ геометрическая кратность всегда не превосходит алгебраическую:
\[ 1 \le \gamma_A(\lambda_0) \le \mu_A(\lambda_0). \]
Матрица $A \in \mathbb{R}^{n \times n}$ диагонализируема тогда и только тогда, когда все корни характеристического
многочлена лежат в поле и для каждого собственного значения геометрическая кратность совпадает с алгебраической:
$\gamma_A(\lambda_i) = \mu_A(\lambda_i)$.
    \end{theorem}

    \section{Линейная независимость собственных векторов для различных собственных значений}

    \begin{theorem}
        Собственные векторы $v_1, v_2, \dots, v_k$, отвечающие попарно различным собственным значениям $\lambda_1,
        \lambda_2, \dots, \lambda_k$ ($\lambda_i \ne \lambda_j$ при $i \ne j$), линейно независимы.
    \end{theorem}
    \begin{proof}
        Проведем строгое доказательство методом математической индукции по числу векторов $k$.

        \textbf{База индукции ($k = 1$):} По определению собственный вектор $v_1 \ne 0$.
        Любая система из одного ненулевого вектора линейно независима ($c_1 v_1 = 0 \implies c_1 = 0$).
        База доказана.

        \textbf{Индукционный переход:} Предположим, что теорема справедлива для любой системы из $k - 1$ собственных
        векторов, отвечающих различным собственным значениям, то есть система $\{v_1, \dots, v_{k-1}\}$ линейно
        независима.

        Докажем независимость системы $\{v_1, \dots, v_k\}$.
        Рассмотрим тривиальную линейную комбинацию:
        \begin{equation} \label{eq:ind1} c_1 v_1 + c_2 v_2 + \dots + c_{k-1} v_{k-1} + c_k v_k = 0.
        \end{equation}
        Применим оператор $A$ к обеим частям равенства \eqref{eq:ind1}.
        В силу линейности оператора $A$ и равенств $A v_i = \lambda_i v_i$:
        \begin{equation} \label{eq:ind2}
            c_1 \lambda_1 v_1 + c_2 \lambda_2 v_2 + \dots + c_{k-1} \lambda_{k-1} v_{k-1} + c_k \lambda_k v_k = 0.
        \end{equation}
        Теперь умножим исходное равенство \eqref{eq:ind1} на скаляр $\lambda_k$:
        \begin{equation} \label{eq:ind3}
            c_1 \lambda_k v_1 + c_2 \lambda_k v_2 + \dots + c_{k-1} \lambda_k v_{k-1} + c_k \lambda_k v_k = 0.
        \end{equation}
        Вычтем равенство \eqref{eq:ind3} из равенства \eqref{eq:ind2}:
        \[ c_1 (\lambda_1 - \lambda_k) v_1 + c_2 (\lambda_2 - \lambda_k) v_2 + \dots + c_{k-1} (\lambda_{k-1} -
        \lambda_k) v_{k-1} + 0 = 0. \]
        Мы получили линейную комбинацию первых $k - 1$ векторов.
        По индукционному предположению система $\{v_1, \dots, v_{k-1}\}$ линейно независима.
        Следовательно, все коэффициенты этой комбинации обязаны равняться нулю:
        \[ c_i (\lambda_i - \lambda_k) = 0 \quad \forall i \in \{1, 2, \dots, k-1\}. \]
        Так как по условию все собственные значения попарно различны, разность $(\lambda_i - \lambda_k) \ne 0$ для всех
        $i < k$.
        Отсюда немедленно следует:
        \[ c_i = 0 \quad \forall i \in \{1, 2, \dots, k-1\}. \]
        Подставляя $c_1 = c_2 = \dots = c_{k-1} = 0$ в исходное равенство \eqref{eq:ind1}, получаем $c_k v_k = 0$.
        А так как $v_k \ne 0$, получаем $c_k = 0$.
        Таким образом, все коэффициенты $c_1 = c_2 = \dots = c_k = 0$, что доказывает линейную независимость системы
        $\{v_1, \dots, v_k\}$.
    \end{proof}

    \section{Диагонализация и возведение матриц в степень}

    \begin{theorem}[Диагонализация через собственный базис]
Пусть матрица $A \in \mathbb{R}^{n \times n}$ имеет $n$ линейно независимых собственных векторов $v_1, \dots, v_n$,
отвечающих собственным значениям $\lambda_1, \dots, \lambda_n$.
Сформируем матрицу $P$, столбцами которой являются собственные векторы, и диагональную матрицу $\Lambda$:
\[
    P = \begin{pmatrix}
            | & & | \\ v_1 & \dots & v_n \\ | & & |
    \end{pmatrix}, \quad \Lambda = \begin{pmatrix}
                                       \lambda_1 & &
                                       0 \\ & \ddots & \\ 0 & & \lambda_n
    \end{pmatrix}.
\]
Тогда:
\[ A P = P \Lambda \iff A = P \Lambda P^{-1}. \]
    \end{theorem}

    \subsection{Вычисление степеней матриц}
    Если матрица диагонализируема, возведение ее в степень $k \in \mathbb{N}$ тривиализуется:
    \[ A^k=(P \Lambda P^{-1})^k = (P \Lambda P^{-1})(P \Lambda P^{-1}) \dots (P \Lambda P^{-1}) = P \Lambda^k P^{-1}. \]
    Возведение диагональной матрицы $\Lambda$ в степень сводится к возведению ее диагональных элементов в степень $k$:
    \[ \Lambda^k = \diag(\lambda_1^k, \lambda_2^k, \dots, \lambda_n^k). \]
    Сложность прямого умножения матриц составляет $\mathcal{O}(k n^3)$, в то время как через диагонализацию она
    составляет $\mathcal{O}(n^3 + n \log k)$!
    Это позволяет мгновенно исследовать асимптотическое поведение динамических систем, цепей Маркова и разностных
    уравнений при $k \to \infty$.

    \section{Вещественная спектральная теорема для симметричных матриц}

    \begin{theorem}[Вещественная спектральная теорема]
Пусть $A \in \mathbb{R}^{n \times n}$ — вещественная симметричная матрица ($A^\top = A$).
Тогда:
\begin{enumerate}
    \item Все собственные значения матрицы $A$ вещественны: $\lambda_i \in \mathbb{R}$;
    \item Собственные векторы, отвечающие различным собственным значениям, взаимно ортогональны;
    \item Существует \textbf{ортонормированный базис} пространства $\mathbb{R}^n$, состоящий из собственных векторов
    матрицы $A$.
\end{enumerate}
Следовательно, матрица $A$ ортогонально диагонализируема:
\[ A = Q \Lambda Q^\top = Q \Lambda Q^{-1}, \]
где $Q \in \mathbb{R}^{n \times n}$ — ортогональная матрица ($Q^\top Q = Q Q^\top = I$), а $\Lambda = \diag(\lambda_1,
\dots, \lambda_n)$.
    \end{theorem}

    \begin{remark}[Спектральное разложение симметричной матрицы]
Записав ортогональную диагонализацию в блочной форме через столбцы $Q = [q_1 \dots q_n]$, получаем \textbf{спектральное
разложение}:
\[ A = \sum_{i=1}^n \lambda_i q_i q_i^\top. \]
Каждое слагаемое $q_i q_i^\top$ является матрицей ортогонального проектора ранга 1 на одномерное собственное
подпространство $\spanop(q_i)$.
    \end{remark}

    \section{Положительно определенные матрицы, Cholesky и QR}

    \begin{definition}[Положительно определенная матрица]
Симметричная матрица $A \in \mathbb{R}^{n \times n}$ называется \textbf{положительно определённой} (обозначается $A
\succ 0$), если соответствующая квадратичная форма строго положительна для любого ненулевого вектора:
\[ x^\top A x > 0 \quad \forall x \in \mathbb{R}^n \setminus \{0\}. \]
    \end{definition}

    \begin{theorem}[Критерии положительной определённости]
Для симметричной матрицы $A = A^\top$ следующие утверждения эквивалентны:
\begin{enumerate}
    \item $A \succ 0$;
    \item Все собственные значения строго положительны: $\lambda_i(A) > 0$ для всех $i = 1, \dots, n$;
    \item \textbf{Критерий Сильвестра:} Все угловые главные миноры строго положительны ($\Delta_k > 0$ для
    $k = 1, \dots, n$);
    \item Существует единственная невырожденная нижнетреугольная матрица $L$ с положительными элементами на
    диагонали такая, что: \[ A = L L^\top \quad \text{(\textbf{Разложение Холецкого})}. \]
\end{enumerate}
    \end{theorem}

    \begin{definition}[QR-decomposition]
\textbf{$QR$-разложением} матрицы $A \in \mathbb{R}^{m \times n}$ ($m \ge n$) называется факторизация:
\[ A = Q R, \]
где $Q \in \mathbb{R}^{m \times n}$ имеет ортонормированные столбцы ($Q^\top Q = I_n$), а $R \in \mathbb{R}^{n \times
n}$ — верхнетреугольная невырожденная матрица с положительными диагональными элементами.
Разложение получается в результате ортогонализации Грама-Шмидта столбцов матрицы $A$.
    \end{definition}

    \section{Сравнительный анализ матричных разложений}

    В таблице~\ref{tab:decompositions} приведен детальный сравнительный обзор основных разложений курса.

    \begin{table}[htbp]
        \centering
        \small
        \begin{tabular}{lp{3.5cm}p{3cm}cp{4cm}}
            \toprule
            \textbf{Разложение} & \textbf{Требования к матрице} & \textbf{Вид сомножителей} & \textbf{Сложность} &
            \textbf{Основное применение} \\ \midrule
            $LU$ & Квадратная $n \times n$, обратима (с перестановками $P$) & $L$ нижнетреугольная, $U$
            верхнетреугольная & $\frac{2}{3} n^3$ & Решение СЛУ $Ax=b$, вычисление определителя и обратной матрицы \\
            \midrule
            Холецкий & Квадратная $n \times n$, $A = A^\top \succ 0$ & $A = L L^\top$, $L$ нижнетреугольная &
            $\frac{1}{3} n^3$ & Оптимизация, МНК, многомерные нормальные распределения (в 2 раза быстрее LU) \\
            \midrule
            $QR$ & Произвольная $m \times n$, $m \ge n$ & $Q$ ортонормированная, $R$ верхнетреугольная
            & $2 n^2 \left(m - \frac{n}{3}\right)$ & Задача МНК, численный алгоритм поиска собственных значений
            (QR-алгоритм) \\ \midrule
            Спектральное & Квадратная $n \times n$, $A = A^\top$
            & $A = Q \Lambda Q^\top$, $Q$ ортогональная, $\Lambda$ диагональная & $\mathcal{O}(n^3)$ & Анализ
            квадратичных форм, PCA, физические колебания, динамические системы \\ \midrule
            $SVD$ & \textbf{Любая} прямоугольная $m \times n$ & $U, V$ ортогональные, $\Sigma$ псевдодиагональная
            & $\mathcal{O}(m n^2)$
            & Низкоранговая аппроксимация, сжатие данных, псевдообращение, анализ главных направлений \\ \bottomrule
        \end{tabular}
        \caption{Сводная классификация фундаментальных матричных разложений}
        \label{tab:decompositions}
    \end{table}

    \newpage

% =========================================================================
% ГЛАВА 4
% =========================================================================